\section{Anlageurteil}
\label{sec:urteil-rdy}

\subsection{Was belegt ist}

Der Umsatz stieg FY2026 um \Pct{\RdyUmsatzWachstum} auf
\MioINR{\RdyUmsatzXXVI}, w\"ahrungsbereinigt fiel er um
\PctBetrag{\RdyOrganisch}. Der Rohertrag fiel, die Rohertragsmarge sank von
\Pct{\RdyRohertragsmargeVJ} auf \Pct{\RdyRohertragsmarge}, das
Betriebsergebnis um \PctBetrag{\RdyBetriebsergebnisWachstum} und der
Jahres\"uberschuss um \PctBetrag{\RdyErgebnisWachstum}. Im Quartal zum
30.~Juni 2026 fiel der Gewinn um \PctBetrag{\RdyQuartalGewinnWachstum} und
der operative Mittelzufluss war negativ. Nordamerika verlor
\MioINR{\RdyNordamerikaRueckgang}, Europa gewann
\MioINR{\RdyEuropaZuwachs}, davon \MioINR{\RdyNrtUmsatz} aus einem Zukauf.
Das Unternehmen h\"alt eine Nettokasse von \MioINR{\RdyNettokasse}, sch\"uttet
\Pct{\RdyAusschuettungsquote} des Gewinns aus und gibt keine Prognose ab. Der
Umsatzbeitrag des Lenalidomid-Gesch\"afts ist in keiner der ausgewerteten
Quellen beziffert. Der Kurs steht \PctBetrag{\RdyAbstandHoch} unter dem
Jahreshoch; \"uber zehn Jahre erreichte der Titel \Pct{\RdyRenditeZehn} pro
Jahr gegen eine Branchenrendite von \Pct{\Huerde}.

\subsection{Was das Urteil tr\"agt}

\bk{\emph{Die Bilanz ist nicht das Problem.} Nettokasse, eine
Eigenkapitalquote von \Pct{\RdyEigenkapitalquote} und eine
Aussch\"uttungsquote unter einem F\"unftel: Dieses Unternehmen ger\"at nicht
in Not. Was immer die n\"achsten Jahre bringen, sie bringen keine
Finanzierungsfrage.}

\bk{\emph{Das Problem ist die Ertragsbasis.} Zwei aufeinanderfolgende
Berichtsperioden zeigen dasselbe Muster: Der Umsatz h\"alt sich \"uber
W\"ahrung und Zukauf, das Ergebnis f\"allt. Das ist die Signatur eines
Gesch\"afts, dessen margenstarker Teil wegbricht und durch margenschw\"acheren
ersetzt wird. Der Rohertrag des gr\"o{\ss}ten Segments fiel um
\Pct{\RdyRohertragGgWachstum}, w\"ahrend sein Umsatz um
\Pct{\RdyGlobalGenericsWachstum} stieg.}

\bk{\emph{Und die entscheidende Gr\"o{\ss}e ist nicht ver\"offentlicht.} Wie
gro{\ss} der wegbrechende Teil ist, sagt das Unternehmen nicht. Es nennt eine
Preissenkung, einen Einmaleffekt von \MioINR{\RdySsaEffekt} und einen
R\"uckgang des Therapiegebiets Onkologie um \MioINR{\RdyOnkologieRueckgang},
stellt zwischen den dreien aber keinen Zusammenhang her. Wer diesen Titel
bewerten will, muss die wichtigste Zahl sch\"atzen. Dieses Dokument tut das
nicht.}

\subsection{Nicht verf\"ugbare Angaben}

Der Umsatz- und Ergebnisbeitrag des Lenalidomid-Gesch\"afts sowie sein
Auslaufpfad: nicht beziffert, in f\"unf Einreichungen \"uber sieben Jahre
gesucht (Abschnitt~\ref{sec:lenalidomid-rdy}).

Eine Prognose des Unternehmens: existiert nicht
(Abschnitt~\ref{sec:prognose-rdy}). Damit fehlt der Ma{\ss}stab, an dem
Teil~\ref{sec:prognose-azn} die Unternehmensleitung von AstraZeneca messen
konnte.

Eine Aufteilung des Quartalsergebnisses nach Segmenten liegt nicht in der
Gliederung des Jahresberichts vor; der Quartalsbericht nennt Gr\"unde im
Flie{\ss}text, keine Segmentzahlen.

Die Sch\"atzungen der indischen Analysten zur Stammaktie sind nicht erhoben.
Erhoben ist der Konsens zum Hinterlegungsschein, und er beruht auf
\Stueck{\RdyKonsensHaeuser} Sch\"atzungen, weniger als die H\"alfte dessen,
was f\"ur AstraZeneca vorliegt.
